\scanpage{61}
\mspage{26}
\diagram{\begin{tikzcd}
\Sigma^{<}(1)\arrow[r,shift left=2,"p_1=\uline{s}"]\arrow[r,shift right=2,"\uline{b}"'] & \Sigma
\end{tikzcd}\qquad
\begin{tikzcd}
\Sigma^{<}(2)\arrow[r,shift left=4,"p_2"]\arrow[r,"\uline{s'}" description]\arrow[r,shift right=4,"\uline{b'}"'] & \Sigma^{<}(1)
\end{tikzcd}}
comme
\[
(39)\qquad\left\{\begin{aligned}
(x\prec y)&\xrightarrow{p_1=\uline{s}}x\\
(x\prec y)&\xrightarrow{\uline{b}}y\\
(x\prec y\prec z)&\xrightarrow{p_2}(x,y)\\
(x\prec y\prec z)&\xrightarrow{\uline{s'}}(\deleted{y}\inserted{x},z)\\
(x<y<z)&\xrightarrow{\uline{b'}}(\deleted{x}\inserted{y},z)
\end{aligned}\right.
\]
\uline{Scholie} En résumé, la situation est décrite
par les données suivantes

$1^{\circ})$ Un espace top. \deleted{(séparé)} $\Sigma$, munie d'une stratifica-
tion (pré)cylindrique --- soit $\Sigma_1$ l'espace des strates de
codim $1$, de sorte que $\Sigma^{<}(1)\xrightarrow{\uline{b}}\Sigma$ est une imm. bordante
\inserted{propre}
\margin{``stratification cylindrique localt ordonnée''}

$2^{\circ})$ Une relation d'ordre totale relative de $\Sigma_1$ sur $\Sigma$
(définie par la donnée, pour tt $x\in\uline{b}(\Sigma_1)$, d'un
ordre total sur l'ens des strates de codim $1$
qui passent par $x$, cette relation d'ordre total étant
soumis à la condition d'être \inserted{donnée par} une partie \uline{fermée}
\inserted{$\Sigma^{*}(2)$} de $\bigl((\Sigma^{<}(1),b)\times_\Sigma(\Sigma^{<}(1),b)\ldots\bigr)$ --- soit
\diagram{\begin{tikzcd}
\Sigma^{<}(2)\arrow[r,shift left=2,"\uline{s'}"]\arrow[r,shift right=2,"\uline{b'}"'] & \Sigma^{<}(1)\arrow[r,"\uline{b}"] & \Sigma
\end{tikzcd}}

$3^{\circ})$ Une application continue
\[
p_1\text{ ou }\uline{s}:\Sigma^{<}(1)\longrightarrow\Sigma
\]
ayant les propriétés suivantes

a) $(\uline{s},\uline{b}):\Sigma^{<}(1)\longrightarrow\Sigma\times\Sigma$ est injectif, et
son image ne
rencontre pas la
diagonale \inserted{$\forall\ u\in\Sigma^{<}(1),\ \uline{s}(u)\ne\uline{b}(u)$}
\mspage{60}
